% Lista de acronimos y siglas.
\chapter*{Acrónimos y siglas}
\addcontentsline{toc}{chapter}{Acrónimos y siglas}
\markboth{Acrónimos y siglas}{Acrónimos y siglas}

\begin{longtable}{>{\bfseries}p{0.18\textwidth} p{0.74\textwidth}}
ABM & Modelos basados en agentes (\textit{Agent-Based Models}). \\[2pt]
AC & Autómata celular (\textit{Cellular Automaton}). \\[2pt]
CLUE-S & \textit{Conversion of Land Use and its Effects at Small regional extent}. \\[2pt]
CONAPO & Consejo Nacional de Población. \\[2pt]
CPU & Unidad central de procesamiento (\textit{Central Processing Unit}). \\[2pt]
EPSG & Registro de sistemas de referencia de coordenadas (\textit{European Petroleum Survey Group}). \\[2pt]
ExG & Índice de exceso de verde (\textit{Excess Green Index}). \\[2pt]
ExR & Índice de exceso de rojo (\textit{Excess Red Index}). \\[2pt]
FLUS & \textit{Future Land Use Simulation}. \\[2pt]
FoM & Figura de Mérito (\textit{Figure of Merit}). \\[2pt]
FRAGSTATS & Programa de análisis de métricas del paisaje (\textit{Fragmentation Statistics}). \\[2pt]
GPU & Unidad de procesamiento gráfico (\textit{Graphics Processing Unit}). \\[2pt]
HSV & Espacio de color matiz-saturación-valor (\textit{Hue, Saturation, Value}). \\[2pt]
IIEc & Instituto de Investigaciones Económicas, UNAM. \\[2pt]
IIMAS & Instituto de Investigaciones en Matemáticas Aplicadas y en Sistemas. \\[2pt]
INEGI & Instituto Nacional de Estadística y Geografía. \\[2pt]
IoU & Índice de solapamiento (\textit{Intersection over Union}). \\[2pt]
IV & Valor de Información (\textit{Information Value}). \\[2pt]
LAB & Espacio de color CIELAB (luminosidad y dos ejes cromáticos). \\[2pt]
LBP & Patrones binarios locales (\textit{Local Binary Patterns}). \\[2pt]
LUCC & Cambio en la cobertura y uso del suelo (\textit{Land Use / Land Cover Change}). \\[2pt]
NDBI & Índice de edificación de diferencia normalizada (\textit{Normalized Difference Built-up Index}). \\[2pt]
NDVI & Índice de vegetación de diferencia normalizada (\textit{Normalized Difference Vegetation Index}). \\[2pt]
NGRDI & Índice verde--rojo normalizado (\textit{Normalized Green--Red Difference Index}). \\[2pt]
NIR & Infrarrojo cercano (\textit{Near Infrared}). \\[2pt]
PCA & Análisis de componentes principales (\textit{Principal Component Analysis}). \\[2pt]
PCIC & Posgrado en Ciencia e Ingeniería de la Computación. \\[2pt]
PMOTDU & Programa Municipal de Ordenamiento Territorial y Desarrollo Urbano. \\[2pt]
RGB & Modelo de color rojo-verde-azul (\textit{Red-Green-Blue}). \\[2pt]
ROC & Característica operativa del receptor (\textit{Receiver Operating Characteristic}). \\[2pt]
SECIHTI & Secretaría de Ciencia, Humanidades, Tecnología e Innovación. \\[2pt]
SEDATU & Secretaría de Desarrollo Agrario, Territorial y Urbano. \\[2pt]
SLEUTH & \textit{Slope, Land use, Exclusion, Urban, Transportation, Hillshade}. \\[2pt]
SVM & Máquina de vectores de soporte (\textit{Support Vector Machine}). \\[2pt]
UML & Lenguaje unificado de modelado (\textit{Unified Modeling Language}). \\[2pt]
UNAM & Universidad Nacional Autónoma de México. \\[2pt]
WGS84 & Sistema geodésico mundial de 1984 (\textit{World Geodetic System 1984}). \\[2pt]
WoE & Pesos de Evidencia (\textit{Weight of Evidence}). \\[2pt]
ZMCM & Zona Metropolitana de la Ciudad de México. \\[2pt]
ZMQ & Zona Metropolitana de Querétaro. \\
\end{longtable}
\cleardoublepage
